\documentclass[12pt]{article}
\NeedsTeXFormat{LaTeX2e}
\usepackage{/home/dennis/src/principia/principia-master/principia} 
%\usepackage{fullpage}
%\usepackage[T1]{fontenc}
%\usepackage[utf8]{inputenc}
%\usepackage{setspace}
%\usepackage{amssymb} 
%\usepackage{amsmath} 
%\usepackage{pifont} 
%\usepackage{graphicx}
\usepackage{hyperref}
\usepackage{color}
%Volume I
%Mathematical logic
%The theory of deduction
%Meta-logical symbols
\newcommand{\ie}{\textit{i}.\textit{e}.\ }
\newcommand{\Ie}{\textit{I}.\textit{e}.\ }
\newcommand{\eg}{\textit{e}.\textit{g}.\ }
\newcommand{\Eg}{\textit{E}.\textit{g}.\ }
\newcommand{\pmsch}[1]{\pmast#1} %Starred chapter
\newcommand{\pmschs}[2]{\pmast#1\text{---}\pmast#2} %Starred chapter
\newcommand{\pmsns}[3]{\pmast#1\pmcdot#2\text{---}\pmcdot#3}%Starred number
\newcommand{\pmpsn}[2]{(\pmast#1\pmcdot#2)} 
\newcommand{\pmpsnn}[3]{(\pmast#1\pmcdot#2\pmcdot#3)} 
\newcommand{\pmsn}[2]{\pmast#1\pmcdot#2} 
\newcommand{\pmnsn}[1]{\text{#1}}
\newcommand{\pmsnn}[3]{\pmast#1\pmcdot#2\pmcdot#3}
\newcommand{\pmsnnn}[4]{\pmast#1\pmcdot#2\pmcdot#3\pmcdot#4}
\newcommand{\pmsnnnn}[5]{\pmast#1\pmcdot#2\pmcdot#3\pmcdot#4\pmcdot#5}
\newcommand{\pmsnnnnn}[6]{\pmast#1\pmcdot#2\pmcdot#3\pmcdot#4\pmcdot#5\pmcdot#6}
\newcommand{\pmsnb}[2]{\boldsymbol{\pmast#1\pmcdot#2}} %Starred number boldface
\newcommand{\pmsnnb}[3]{\boldsymbol{\pmast#1\pmcdot#2\pmcdot#3}}
\newcommand{\pmsnnnb}[4]{\boldsymbol{\pmast#1\pmcdot#2\pmcdot#3\pmcdot#4}}
\newcommand{\pmsnnnnb}[5]{\boldsymbol{\pmast#1\pmcdot#2\pmcdot#3\pmcdot#4\pmcdot#5}}
\newcommand{\pmsnnnnnb}[6]{\boldsymbol{\pmast#1\pmcdot#2\pmcdot#3\pmcdot#4\pmcdot#5\pmcdot#6}}
\newcommand{\pmfd}{\begin{center} \rule{5cm}{.5pt} \end{center}} %Dividing line between introductory remarks in a starred number and the formal deductions.
\newcommand{\pmdem}{\textit{Dem}.} %This notation begins a proof.
\newcommand{\pmdemi}{\indent \pmdem} %This idents the notation that begins a proof.
\newcommand{\pmhp}{\text{Hp}} %This typesets Hp (short for antecedent), which occurs at the beginning of a proof.
\newcommand{\pmprop}{\text{Prop}} %This occurs at the end of a proof.
\newcommand{\pmithm}{\pmimp\;\pmthm} %This occurs when a meta-theoretic implication is asserted.
\newcommand{\pmbr}[1]{\bigg \lbrack \normalsize #1 \bigg \rbrack} %These are larger brackets for substitution.
\newcommand{\pmsub}[2]{\bigg \lbrack \small \begin{array}{c} #1 \\ \hline #2 \end{array} \bigg \rbrack} %This is the substitution command.
\newcommand{\pmsubb}[4]{\bigg \lbrack \small \begin{array}{c c} #1, & #3 \\ \hline #2, & #4 \end{array}  \bigg \rbrack} %This is the substitution command.
\newcommand{\pmsubbb}[6]{\bigg \lbrack \small \begin{array}{c c c} #1, & #3, & #5 \\ \hline #2, & #4, & #6 \end{array}  \bigg \rbrack} %This is the substitution command.
\newcommand{\pmsubbbb}[8]{\bigg \lbrack \small \begin{array}{c c c c} #1, & #3, & #5, & #7 \\ \hline #2, & #4, & #6, & #8 \end{array}  \bigg \rbrack} %This is the substitution command.
\newcommand{\pmSub}[3]{\bigg \lbrack \normalsize #1 \text{ } \small \begin{array}{c} #2 \\ \hline #3 \end{array}  \bigg \rbrack} %This is the substitution command.
\newcommand{\pmSubb}[5]{\bigg \lbrack \normalsize #1 \text{ } \small \begin{array}{c c} #2, & #4 \\ \hline #3, & #5 \end{array}  \bigg \rbrack} %This is the substitution command.
\newcommand{\pmSubbb}[7]{\bigg \lbrack \normalsize #1 \text{ } \small \begin{array}{c c c} #2, & #4, & #6 \\ \hline #3, & #5, & #7 \end{array}  \bigg \rbrack} %This is the substitution command.
\newcommand{\pmSubbbb}[9]{\bigg \lbrack \normalsize #1 \text{ } \small \begin{array}{c c c c} #2, & #4, & #6, & #8 \\ \hline #3, & #5, & #7, & #9 \end{array} \bigg \rbrack} %This is the substitution command.
\newcommand{\pmsUb}[2]{\small \begin{array}{c} #1 \\ \hline #2 \end{array}} %This is the substitution command.
\newcommand{\pmsUbb}[4]{\small \begin{array}{c c} #1, & #3 \\ \hline #2, & #4 \end{array}} %This is the substitution command.
\newcommand{\pmsUbbb}[6]{\small \begin{array}{c c c} #1, & #3, & #5 \\ \hline #2, & #4, & #6 \end{array}} %This is the substitution command.
\newcommand{\pmsUbbbb}[8]{\small \begin{array}{c c c c} #1, & #3, & #5, & #7 \\ \hline #2, & #4, & #6, & #8 \end{array}} %This is the substitution command.
\newcommand{\pmSUb}[3]{\normalsize #1 \text{ } \small \begin{array}{c} #2 \\ \hline #3 \end{array}} %This is the substitution command.
\newcommand{\pmSUbb}[5]{\normalsize #1 \text{ } \small \begin{array}{c c} #2, & #4 \\ \hline #3, & #5 \end{array}} %This is the substitution command.
\newcommand{\pmSUbbb}[7]{\normalsize #1 \text{ } \small \begin{array}{c c c} #2, & #4, & #6 \\ \hline #3, & #5, & #7 \end{array}} %This is the substitution command.
\newcommand{\pmSUbbbb}[9]{\normalsize #1 \text{ } \small \begin{array}{c c c c} #2, & #4, & #6, & #8 \\ \hline #3, & #5, & #7, & #9 \end{array}} %This is the substitution command.
\newcommand{\pmthm}{\mathpunct{\text{\scalebox{.5}[1]{$\boldsymbol\vdash$}}}} %This is the theorem sign.
\newcommand{\pmast}{\text{\resizebox{!}{.75\height}{\ding{107}}}} %This is the sign introducing a theorem number.
\newcommand{\pmcdot}{\text{\raisebox{.05cm}{$\boldsymbol\cdot$}}} %This is a sign introducing a theorem sub-number.
\newcommand{\pmiddf}{\mathbin{=}}
\newcommand{\pmdf}{\quad \text{Df}}
\newcommand{\pmDf}{\text{Df}}
\newcommand{\pmpp}{\quad \text{Pp}}

%Square dots for scope, defined for up to six dots
\newcommand{\pmd}{\hbox{\rule{.3ex}{.3ex}}} %single square dot
\newcommand{\pmdd}{\overset{\pmd}{\pmd}} %vertically aligned pair of square dots
\newcommand{\pmdot}{\mathinner{\pmd}}
 
 
